% ====================================================================
% Hello Buddy — Product Design Report (PDR)
% Format: Per Product_Design_Report_Documentation_Manual
% ====================================================================
\documentclass[12pt,a4paper]{article}

\usepackage[margin=1in]{geometry}
\usepackage[T1]{fontenc}
\usepackage{mathptmx}
\usepackage{array}
\usepackage{longtable}
\usepackage{enumitem}
\usepackage{tabularx}
\usepackage{ragged2e}
\usepackage{graphicx}
\usepackage{fancyhdr}
\usepackage{hyperref}
\usepackage{xcolor}
\usepackage{tcolorbox}
\usepackage{booktabs}
\usepackage{caption}
\usepackage{setspace}

\setlength{\parindent}{0pt}
\setlength{\parskip}{6pt}
\renewcommand{\arraystretch}{1.35}
\setstretch{1.15}

% Column types
\newcolumntype{L}[1]{>{\RaggedRight\arraybackslash}p{#1}}
\newcolumntype{C}[1]{>{\centering\arraybackslash}p{#1}}
\newcolumntype{Y}{>{\RaggedRight\arraybackslash}X}

% Traceability tag command
\newtcolorbox{tracebox}{
    colback=gray!10, colframe=gray!50, boxrule=0.5pt,
    left=6pt, right=6pt, top=3pt, bottom=3pt,
    fontupper=\small\itshape
}

% Header/Footer
\pagestyle{fancy}
\fancyhf{}
\lhead{\small School of Electronics Engineering --- Electronic Product Development}
\rhead{\small Hello Buddy --- PDR}
\cfoot{\small Page \thepage}
\renewcommand{\headrulewidth}{0.4pt}

\begin{document}

% ====================================================================
%                        COVER PAGE
% ====================================================================
\begin{titlepage}
\begin{center}

\vspace*{1cm}
{\large\bfseries KALINGA INSTITUTE OF INDUSTRIAL TECHNOLOGY}\\[0.2cm]
{\normalsize School of Electronics Engineering}\\[0.2cm]
{\normalsize Bhubaneswar, Odisha}

\vspace{1.5cm}
{\LARGE\bfseries PRODUCT DESIGN REPORT}\\[0.5cm]
{\Large\bfseries Hello Buddy}\\[0.2cm]
{\large Desktop Companion and Study Assistant Robot}

\vspace{1.5cm}
\begin{tabular}{|L{0.35\textwidth}|L{0.50\textwidth}|}
\hline
\textbf{Course} & Electronic Product Development (Sessional) \\ \hline
\textbf{Document Type} & Product Design Report (PDR) \\ \hline
\textbf{Prerequisite} & Requirements Specification Document (RSD) --- Submitted \\ \hline
\textbf{Group Number} & 09 \\ \hline
\textbf{Document Number} & G09\_HelloBuddy\_PDR\_v1 \\ \hline
\textbf{Revision} & 1.0 \\ \hline
\textbf{Date} & September 18, 2026 \\ \hline
\end{tabular}

\vspace{1cm}
\textbf{Team Members}

\begin{tabular}{|L{0.40\textwidth}|L{0.25\textwidth}|L{0.20\textwidth}|}
\hline
\textbf{Name} & \textbf{Roll Number} & \textbf{Role} \\ \hline
Ziaur Rahaman & 2404154 & Firmware \& App Dev \\ \hline
Aryan Raj & 2404086 & Firmware Development \\ \hline
Ayush Dutta & 2404088 & CAD Design \\ \hline
Shamil Minhaj & 2504805 & Circuit Design \\ \hline
\end{tabular}

\vspace{1cm}
\textbf{Faculty-in-Charge:} Dr.\ Deepak Kumar Rout

\vfill
{\small Reference Frameworks: IEEE, ISO 9001/ISO-IEC-IEEE 15288, IPC-2221,
IEC 62368-1, IoT-A Design Methodology}

\end{center}
\end{titlepage}

% ====================================================================
%                   DOCUMENT CONTROL SHEET
% ====================================================================
\newpage
\section*{Document Control Sheet}

\begin{center}
\begin{tabular}{|L{0.30\textwidth}|L{0.58\textwidth}|}
\hline
\textbf{Document Title} & Hello Buddy --- Product Design Report \\ \hline
\textbf{Document Number} & G09\_HelloBuddy\_PDR\_v1 \\ \hline
\textbf{Revision} & 1.0 \\ \hline
\textbf{Prepared By} & Ziaur Rahaman, Aryan Raj, Ayush Dutta, Shamil Minhaj \\ \hline
\textbf{Reviewed By} & Group 09 (Peer Review) \\ \hline
\textbf{Approved By} & Dr.\ Deepak Kumar Rout (Faculty-in-Charge) \\ \hline
\textbf{Approval Status} & Submitted for Review \\ \hline
\textbf{RTM Reference} & See Section~1 (Requirements Traceability Matrix) \\ \hline
\end{tabular}
\end{center}

\vspace{0.5cm}

\subsection*{Revision History}

\begin{center}
\begin{tabular}{|C{0.10\textwidth}|L{0.18\textwidth}|L{0.40\textwidth}|L{0.18\textwidth}|}
\hline
\textbf{Rev.} & \textbf{Date} & \textbf{Description} & \textbf{Author} \\ \hline
0.1 & September 10, 2026 & Initial draft --- hardware and software block diagrams & Ziaur Rahaman \\ \hline
0.2 & September 14, 2026 & Added IoT views, safety checklists, and BOM & All Members \\ \hline
1.0 & September 18, 2026 & Final submission --- all sections complete, peer reviewed & All Members \\ \hline
\end{tabular}
\end{center}

% ====================================================================
%              ABSTRACT, ABBREVIATIONS, TOC
% ====================================================================
\newpage
\section*{Abstract}

This Product Design Report (PDR) presents the complete design of
\textbf{Hello Buddy}, a stationary desktop companion and study-assistant
robot built around the ESP32-S3 N16R8 microcontroller. The design
realises the functional and non-functional requirements approved in the
Requirements Specification Document (RSD, G09\_HelloBuddy\_RSD\_v1).

Hello Buddy features a GC9A01 round TFT display for animated facial
expressions, an SSD1306 OLED for task and session status, four SG90
servos for subtle body movement, an INMP441 I2S microphone with
wake-word activation, VL53L0X ToF proximity sensing, DFPlayer Mini
audio output, and dual TTP223 capacitive touch sensors. Natural language
processing is handled via cloud LLM APIs (Groq, DeepSeek, Gemini, or a
locally hosted model such as Ollama) through a FastAPI backend server
running Whisper STT. Task persistence and session analytics are managed
through Firebase Realtime Database, accessible from both the ESP32-S3
firmware and a companion Flutter mobile application.

The report follows the structure mandated by the PDR Documentation
Manual and is aligned with IEEE, ISO 9001/ISO-IEC-IEEE 15288, IPC-2221,
IEC 62368-1, and the IoT-A design methodology.

\section*{List of Abbreviations}

\begin{center}
\begin{tabular}{|L{0.18\textwidth}|L{0.70\textwidth}|}
\hline
\textbf{Abbreviation} & \textbf{Full Form} \\ \hline
API & Application Programming Interface \\ \hline
BLE & Bluetooth Low Energy \\ \hline
BOM & Bill of Materials \\ \hline
CRUD & Create, Read, Update, Delete \\ \hline
FPS & Frames Per Second \\ \hline
GPIO & General-Purpose Input/Output \\ \hline
I2C & Inter-Integrated Circuit \\ \hline
I2S & Inter-IC Sound \\ \hline
LLM & Large Language Model \\ \hline
OLED & Organic Light-Emitting Diode \\ \hline
OTA & Over-The-Air (update) \\ \hline
PCB & Printed Circuit Board \\ \hline
PDR & Product Design Report \\ \hline
PSRAM & Pseudo Static Random-Access Memory \\ \hline
PWM & Pulse-Width Modulation \\ \hline
REST & Representational State Transfer \\ \hline
RSD & Requirements Specification Document \\ \hline
RTM & Requirements Traceability Matrix \\ \hline
SPI & Serial Peripheral Interface \\ \hline
STT & Speech-To-Text \\ \hline
TFT & Thin-Film Transistor \\ \hline
ToF & Time of Flight \\ \hline
TTS & Text-To-Speech \\ \hline
UART & Universal Asynchronous Receiver-Transmitter \\ \hline
UBEC & Universal Battery Eliminator Circuit \\ \hline
\end{tabular}
\end{center}

\tableofcontents
\listoffigures
\listoftables

% ====================================================================
%     SECTION 1: REQUIREMENTS TRACEABILITY MATRIX
% ====================================================================
\newpage
\section{Requirements Traceability Matrix}

\begin{tracebox}
Traceability --- This section provides forward and backward traceability
for all RSD requirements to design elements and test cases.
\end{tracebox}

\vspace{0.3cm}

Table~\ref{tab:rtm} provides the master Requirements Traceability Matrix
linking each RSD requirement to its design realisation in this PDR and
the corresponding verification method and test case.

\begin{center}
\scriptsize
\captionof{table}{Requirements Traceability Matrix (RTM)}
\label{tab:rtm}
\begin{tabular}{|C{0.07\textwidth}|L{0.22\textwidth}|C{0.04\textwidth}|L{0.22\textwidth}|L{0.13\textwidth}|C{0.06\textwidth}|C{0.08\textwidth}|}
\hline
\textbf{Req.~ID} & \textbf{Requirement Summary} & \textbf{Pri.} &
\textbf{Design Element / PDR Section} & \textbf{Verification} &
\textbf{Test ID} & \textbf{Status} \\ \hline
HW-001 & GC9A01 display $\geq$10\,FPS expressions & H &
\S3 HW Block, \S5 Circuit & Performance Test & TC-01 & Met \\ \hline
HW-002 & SSD1306 OLED task/timer display & H &
\S3 HW Block, \S5 Circuit & Functional Test & TC-02 & Met \\ \hline
HW-003 & INMP441 mic, wake-word only & H &
\S3 HW Block, \S5 Circuit & Inspection & TC-03 & Met \\ \hline
HW-004 & VL53L0X presence 10--60\,cm & H &
\S3 HW Block, \S5 Circuit & Lab Test & TC-04 & Met \\ \hline
HW-005 & DFPlayer MP3 $<$300\,ms latency & H &
\S3 HW Block, \S5 Circuit & Timing Test & TC-05 & Met \\ \hline
HW-006 & 4x SG90 servos jitter-free & H &
\S3 HW Block, \S5 Circuit & Stress Test & TC-06 & Met \\ \hline
HW-007 & Battery $\geq$6\,h operation & M &
\S3 HW Block, \S15 BOM & Duration Test & TC-07 & Partial \\ \hline
HW-008 & 3.3\,V/5\,V only, UBEC protection & H &
\S3 HW Block, \S14 Safety & Inspection & TC-08 & Met \\ \hline
FW-001 & FreeRTOS concurrent tasks & H &
\S4 SW Block & Stress Test & TC-09 & Met \\ \hline
FW-002 & Wake-word ``Hey Buddy'' & H &
\S4 SW Block & Functional Test & TC-10 & Met \\ \hline
FW-003 & Voice pipeline $\leq$3\,s end-to-end & H &
\S4 SW Block, \S10 Service & Timing Test & TC-11 & Met \\ \hline
FW-004 & 8 emotion states on GC9A01 & H &
\S4 SW Block, \S12 Functional & Demonstration & TC-12 & Met \\ \hline
FW-005 & OLED live session data & H &
\S4 SW Block & Functional Test & TC-13 & Met \\ \hline
FW-006 & Touch reaction $<$200\,ms & H &
\S4 SW Block & Timing Test & TC-14 & Met \\ \hline
FW-007 & Absence detection $>$2\,min & M &
\S4 SW Block & Functional Test & TC-15 & Met \\ \hline
FW-008 & HTTP/WebSocket server port 80 & H &
\S4 SW Block, \S13 Operational & Integration Test & TC-16 & Met \\ \hline
API-001 & Whisper STT $\leq$1.5\,s & H &
\S4 SW Block, \S10 Service & Timing Test & TC-17 & Met \\ \hline
API-002 & LLM API JSON action plan & H &
\S4 SW Block, \S10 Service & Integration Test & TC-18 & Met \\ \hline
API-003 & Firebase sync all clients & H &
\S4 SW Block, \S9 Info Model & Integration Test & TC-19 & Met \\ \hline
API-004 & Task CRUD endpoints & H &
\S10 Service & Functional Test & TC-20 & Met \\ \hline
API-005 & Session start/end logging & H &
\S10 Service & Functional Test & TC-21 & Met \\ \hline
API-006 & API keys server-side only & H &
\S14 Safety Checklist & Code Review & TC-22 & Met \\ \hline
API-007 & Spotify control via Spotipy & L &
\S10 Service, \S13 Operational & Functional Test & TC-23 & Partial \\ \hline
APP-001 & Dashboard: timer, task, status & H &
\S4 SW Block, \S12 Functional & UI Test & TC-24 & Met \\ \hline
APP-002 & Task CRUD sync $\leq$1\,s & H &
\S4 SW Block, \S9 Info Model & Functional Test & TC-25 & Met \\ \hline
APP-003 & Pomodoro 25/5 session control & H &
\S4 SW Block & Functional Test & TC-26 & Met \\ \hline
APP-004 & Manual expression/movement control & M &
\S12 Functional & UI Demo & TC-27 & Met \\ \hline
APP-005 & Stats: hours, completion, streaks & M &
\S9 Info Model, \S12 Functional & Functional Test & TC-28 & Met \\ \hline
\end{tabular}
\end{center}

% ====================================================================
%     SECTION 2: INTRODUCTION
% ====================================================================
\newpage
\section{Introduction}

\begin{tracebox}
Traceability --- Requirement(s) addressed: --- (contextual section;
bridges the approved RSD to this design document).
\end{tracebox}

\vspace{0.3cm}

\subsection{Background and Motivation}

Passive digital tools such as phone reminders and desktop widgets are
easily dismissed and fail to create a sense of accountability. Existing
physically expressive robots are either toy-grade or prohibitively
expensive (e.g., Loona at Rs.\,25{,}000+, ElectronBot at Rs.\,8{,}000+).
Hello Buddy addresses this gap as an affordable, open-source, stationary
desktop companion that combines physical emotional expression with
cloud-backed study assistance, targeted at students.

\subsection{Approved RSD Requirements (Recap)}

The following requirements were approved in the RSD
(G09\_HelloBuddy\_RSD\_v1):

\begin{itemize}[leftmargin=*]
\item \textbf{Functional:} Wake-word detection (``Hey Buddy''),
      proximity-based auto-wake, touch reactions, animated facial
      expressions on GC9A01, OLED task/timer display, voice-driven task
      management via LLM API, Pomodoro study sessions, Spotify playback
      control.
\item \textbf{Non-Functional:} API response $<$2\,s, battery
      $\geq$6\,h, concurrent FreeRTOS tasks without jitter, wake-word-only
      microphone activation (privacy-by-design), 3.3\,V/5\,V circuits
      only.
\item \textbf{External Interfaces:} LLM API (Groq/DeepSeek/Gemini/local),
      Firebase Realtime DB, Spotify Web API, Flutter app via
      HTTP/WebSocket.
\item \textbf{Design Constraints:} ESP32-S3 Arduino IDE/FreeRTOS; fits
      inside EMOtto STL body; stationary desktop design; no mains power.
\end{itemize}

\subsection{Scope of This PDR}

This PDR presents the hardware block diagram, software architecture,
circuit schematics, IoT design-methodology views (Process, Domain,
Information, Service, IoT-Level, Functional, Operational), safety
checklists, BOM, and testing results for Hello Buddy, realising all
approved RSD requirements.

% ====================================================================
%     SECTION 3: HARDWARE BLOCK DIAGRAM
% ====================================================================
\newpage
\section{Hardware Block Diagram}

\begin{tracebox}
Traceability --- Requirement(s) addressed: HW-001 through HW-008.
\end{tracebox}

\vspace{0.3cm}

Figure~\ref{fig:hwblock} shows the top-level hardware architecture of
Hello Buddy. The system is centred on the ESP32-S3 N16R8 microcontroller,
with six functional blocks: Power, Sensing, Processing, Display,
Audio, and Actuation.

\vspace{0.5cm}
\begin{center}
\textit{[Insert Hardware Block Diagram --- use KiCad/draw.io/Fritzing.
Show: Power block (LiPo $\rightarrow$ UBEC $\rightarrow$ 5\,V rail),
Processing (ESP32-S3), Display (GC9A01 via SPI + SSD1306 via I2C),
Audio (INMP441 via I2S + DFPlayer via UART + Speaker), Sensing
(VL53L0X via I2C + TTP223 $\times$2 via GPIO), Actuation
(SG90 $\times$4 via PWM). Use dashed lines for power, solid arrows
for signal/data.]}
\end{center}
\vspace{0.3cm}
\captionof{figure}{Hello Buddy --- Top-level hardware block diagram}
\label{fig:hwblock}

\vspace{0.5cm}

\subsection{Block Description Table}

\begin{center}
\captionof{table}{Hardware block descriptions and component selection}
\label{tab:hwblocks}
\scriptsize
\begin{tabular}{|L{0.15\textwidth}|L{0.30\textwidth}|L{0.22\textwidth}|L{0.20\textwidth}|}
\hline
\textbf{Block} & \textbf{Function} & \textbf{Component} & \textbf{Interface} \\ \hline
Processing & Central controller; runs FreeRTOS firmware, WiFi,
             all peripheral drivers & ESP32-S3 N16R8 (16\,MB flash,
             8\,MB PSRAM) & --- \\ \hline
Face Display & Animated facial expressions (8+ emotions) &
              GC9A01 1.28'' Round TFT (240$\times$240) & SPI (GPIO 5/17/16/18/23) \\ \hline
Status Display & Task name, timer, Pomodoro count &
                SSD1306 OLED 128$\times$64 & I2C (GPIO 39/40) \\ \hline
Audio Input & Voice capture after wake-word trigger &
             INMP441 I2S MEMS Microphone & I2S (GPIO 41/42/2) \\ \hline
Audio Output & MP3 playback, TTS responses, sound effects &
              DFPlayer Mini + 8\,$\Omega$ Speaker & UART (GPIO 32/4),
              Busy: GPIO 34 \\ \hline
Proximity & User presence detection (10--60\,cm) &
           VL53L0X ToF Sensor & I2C (GPIO 22/21) \\ \hline
Touch & Capacitive touch reaction (2 zones) &
       TTP223 $\times$2 & GPIO 12, GPIO 33 \\ \hline
Actuation & Subtle body tilt and reaction movements &
           SG90 Micro-Servo $\times$4 & PWM (hardware pins) \\ \hline
Power & Battery and regulation &
       7.4\,V 1200\,mAh LiPo + UBEC 5\,V/3\,A & --- \\ \hline
\end{tabular}
\end{center}

\subsection{Power Budget}

\begin{center}
\captionof{table}{Power budget estimation}
\label{tab:powerbudget}
\scriptsize
\begin{tabular}{|L{0.20\textwidth}|C{0.10\textwidth}|C{0.12\textwidth}|C{0.12\textwidth}|C{0.12\textwidth}|C{0.12\textwidth}|}
\hline
\textbf{Block} & \textbf{Voltage} & \textbf{Current (typ.)} &
\textbf{Power (typ.)} & \textbf{Duty Cycle} & \textbf{Avg.\ Power} \\ \hline
ESP32-S3 (WiFi active) & 3.3\,V & 240\,mA & 792\,mW & 100\% & 792\,mW \\ \hline
GC9A01 Display & 3.3\,V & 25\,mA & 83\,mW & 100\% & 83\,mW \\ \hline
SSD1306 OLED & 3.3\,V & 20\,mA & 66\,mW & 100\% & 66\,mW \\ \hline
INMP441 Mic & 3.3\,V & 1.4\,mA & 5\,mW & 10\% & 0.5\,mW \\ \hline
DFPlayer Mini & 5\,V & 200\,mA & 1000\,mW & 20\% & 200\,mW \\ \hline
VL53L0X & 3.3\,V & 19\,mA & 63\,mW & 100\% & 63\,mW \\ \hline
TTP223 $\times$2 & 3.3\,V & 8\,$\mu$A & $<$1\,mW & 100\% & $<$1\,mW \\ \hline
SG90 $\times$4 (moving) & 5\,V & 800\,mA & 4000\,mW & 5\% & 200\,mW \\ \hline
\textbf{Total} & & & & & \textbf{$\sim$1.4\,W} \\ \hline
\end{tabular}
\end{center}

Estimated battery life with 7.4\,V 1200\,mAh ($\approx$8.88\,Wh) at
$\sim$1.4\,W average draw: $\approx$6.3 hours, satisfying requirement
HW-007.

% ====================================================================
%     SECTION 4: SOFTWARE BLOCK DIAGRAM
% ====================================================================
\newpage
\section{Software Block Diagram}

\begin{tracebox}
Traceability --- Requirement(s) addressed: FW-001 through FW-008,
API-001 through API-007, APP-001 through APP-005.
\end{tracebox}

\vspace{0.3cm}

Figure~\ref{fig:swblock} shows the layered software architecture
spanning the ESP32-S3 firmware, the FastAPI backend server, and the
Flutter companion application.

\vspace{0.5cm}
\begin{center}
\textit{[Insert Software Block Diagram --- use draw.io. Layers:
HAL (GPIO/SPI/I2C/I2S/UART/PWM drivers), FreeRTOS Scheduler,
Middleware (WiFi, HTTP/WebSocket, Firebase Client, ArduinoJSON),
Application (Wake Word, Expression Engine, Servo Controller, OLED
Manager, DFPlayer Controller, Sensor Poller, API Bridge).
Separate column for FastAPI Server (Whisper STT, LLM Router,
Task CRUD, Session Manager, Spotify Controller) and Flutter App
(Dashboard, Task Manager, Pomodoro, Control Panel, Stats).
Firebase Realtime DB shared between all three.]}
\end{center}
\vspace{0.3cm}
\captionof{figure}{Hello Buddy --- Layered software architecture}
\label{fig:swblock}

\vspace{0.5cm}

\subsection{Firmware Architecture (ESP32-S3)}

\begin{center}
\captionof{table}{FreeRTOS task allocation}
\label{tab:rtos}
\scriptsize
\begin{tabular}{|L{0.18\textwidth}|C{0.08\textwidth}|C{0.08\textwidth}|L{0.30\textwidth}|L{0.22\textwidth}|}
\hline
\textbf{Task Name} & \textbf{Core} & \textbf{Priority} &
\textbf{Responsibility} & \textbf{Req.\ Traced} \\ \hline
DisplayTask & 1 & 3 & GC9A01 face rendering at $\geq$10\,FPS;
                       OLED session data update every 1\,s &
                       HW-001, HW-002, FW-004, FW-005 \\ \hline
AudioTask & 1 & 2 & DFPlayer MP3 playback, async callback
                     flag management & HW-005, FW-006 \\ \hline
ServoTask & 0 & 4 & SG90 PWM control for body-tilt routines;
                     hardware PWM pins only & HW-006 \\ \hline
SensorTask & 0 & 3 & VL53L0X polling (50\,ms), TTP223 read,
                      presence/absence detection &
                      HW-004, FW-006, FW-007 \\ \hline
WakeWordTask & 0 & 5 & MicroWakeWord ``Hey Buddy'' detection;
                        INMP441 I2S capture on trigger &
                        HW-003, FW-002 \\ \hline
NetworkTask & 0 & 2 & WiFi management, HTTP/WebSocket server,
                       API bridge to FastAPI backend &
                       FW-003, FW-008 \\ \hline
\end{tabular}
\end{center}

\subsection{Main Control Loop --- State Machine}

\vspace{0.3cm}
\begin{center}
\textit{[Insert State Machine / Flowchart --- states: IDLE,
LISTENING, PROCESSING, EXECUTING, STUDY\_SESSION, BREAK, SLEEP.
Transitions: wake-word $\rightarrow$ LISTENING, API response
$\rightarrow$ EXECUTING, timeout $\rightarrow$ SLEEP, touch
$\rightarrow$ reaction, Pomodoro timer $\rightarrow$ BREAK.]}
\end{center}
\vspace{0.3cm}
\captionof{figure}{Firmware main control loop state machine}
\label{fig:statemachine}

\subsection{Backend Server (FastAPI)}

The FastAPI server runs on a PC/laptop on the same WiFi network and
provides the following services:

\begin{center}
\captionof{table}{FastAPI backend endpoints}
\label{tab:apiendpoints}
\scriptsize
\begin{tabular}{|L{0.25\textwidth}|L{0.12\textwidth}|L{0.50\textwidth}|}
\hline
\textbf{Endpoint} & \textbf{Method} & \textbf{Description} \\ \hline
\texttt{/stt} & POST & Receives audio from ESP32-S3, runs Whisper STT,
                        returns text transcript (API-001) \\ \hline
\texttt{/llm} & POST & Forwards transcript to configured LLM API
                        (Groq/DeepSeek/Gemini/Ollama), returns JSON
                        action plan (API-002) \\ \hline
\texttt{/tasks} & GET/POST & Retrieve all tasks / add new task; synced
                             to Firebase (API-003, API-004) \\ \hline
\texttt{/tasks/\{id\}} & PATCH/DELETE & Update or delete task
                                        (API-004) \\ \hline
\texttt{/session/start} & POST & Log study session start with subject
                                 to Firebase (API-005) \\ \hline
\texttt{/session/end} & POST & Log session end with duration and
                               Pomodoro count (API-005) \\ \hline
\texttt{/spotify/play} & POST & Start Spotify playlist via Spotipy
                                (API-007) \\ \hline
\texttt{/spotify/pause} & POST & Pause Spotify playback
                                 (API-007) \\ \hline
\end{tabular}
\end{center}

\subsection{Flutter Companion Application}

The Flutter app provides the user-facing interface for task management,
study sessions, and Hello Buddy control. It communicates with Firebase
Realtime Database for data persistence and optionally connects directly
to the ESP32-S3 HTTP server for low-latency commands.

Key screens: Home Dashboard (APP-001), Task Manager (APP-002),
Pomodoro Session (APP-003), Hello Buddy Control Panel (APP-004),
Statistics and Mood Log (APP-005).

% ====================================================================
%     SECTION 5: CIRCUIT DIAGRAMS
% ====================================================================
\newpage
\section{Circuit Diagrams}

\begin{tracebox}
Traceability --- Requirement(s) addressed: HW-001 through HW-008.
\end{tracebox}

\vspace{0.3cm}

\textbf{Drafting Standards Applied:} All schematics in this section
comply with the standards listed in Table~\ref{tab:drafting}.

\begin{center}
\captionof{table}{Circuit-diagram drafting standards applied}
\label{tab:drafting}
\scriptsize
\begin{tabular}{|L{0.30\textwidth}|L{0.58\textwidth}|}
\hline
\textbf{Standard} & \textbf{What It Governs} \\ \hline
IEC 60617 / ISO 14617 & Graphical symbols for components (resistor,
                        capacitor, transistor, etc.). No mixed
                        US/ANSI and IEC symbols. \\ \hline
IEC 61082-1 & Diagram layout, signal-flow direction (left-to-right /
              top-to-bottom), sheet cross-references. \\ \hline
ISO 5457 & Drawing sheet size (A4) and border/frame layout. \\ \hline
ISO 7200 & Title-block data fields on every schematic sheet. \\ \hline
IEC 61346-2 / IEEE Std 200 & Reference designators (R1, C3, U2)
                              consistent across schematic, component
                              list, and BOM. \\ \hline
ISO 128 & Line conventions: continuous = electrical connection,
          dashed = mechanical/enclosure outline. \\ \hline
\end{tabular}
\end{center}

\subsection{Sub-Circuit: Power Supply and Regulation}

\begin{center}
\textit{[Insert schematic: 7.4\,V LiPo $\rightarrow$ Switch
$\rightarrow$ UBEC (5\,V 3\,A output) $\rightarrow$ 5\,V rail.
3.3\,V derived from ESP32-S3 on-board LDO. Include reverse-polarity
protection diode, decoupling capacitors (100\,$\mu$F + 100\,nF)
at UBEC output. Title block per ISO 7200.]}
\end{center}
\captionof{figure}{Power supply and regulation circuit}
\label{fig:ckt_power}

\subsection{Sub-Circuit: ESP32-S3 Minimum System}

\begin{center}
\textit{[Insert schematic: ESP32-S3 N16R8 module with all GPIO
connections labelled, USB-C programming port, EN/BOOT buttons,
decoupling capacitors at 3V3 pin. Title block per ISO 7200.]}
\end{center}
\captionof{figure}{ESP32-S3 N16R8 minimum system circuit}
\label{fig:ckt_mcu}

\subsection{Sub-Circuit: Display Connections}

\begin{center}
\textit{[Insert schematic: GC9A01 SPI connection (CS=GPIO5, DC=GPIO17,
RST=GPIO16, SCL=GPIO18, SDA=GPIO23, VCC=3.3V, GND). SSD1306 I2C
connection (SCL=GPIO39, SDA=GPIO40, VCC=3.3V, GND, ADDR=0x3C).
Title block per ISO 7200.]}
\end{center}
\captionof{figure}{GC9A01 and SSD1306 display connections}
\label{fig:ckt_display}

\subsection{Sub-Circuit: Audio (INMP441 + DFPlayer Mini)}

\begin{center}
\textit{[Insert schematic: INMP441 I2S (SCK=GPIO41, WS=GPIO42,
SD=GPIO2, VDD=3.3V, GND, L/R=GND for left channel). DFPlayer Mini
UART (RX=GPIO32 via 1\,k$\Omega$ resistor, TX=GPIO4, Busy=GPIO34,
VCC=5V, SPK1/SPK2 to 8$\Omega$ speaker). Title block per ISO 7200.]}
\end{center}
\captionof{figure}{Audio input (INMP441) and output (DFPlayer Mini) circuit}
\label{fig:ckt_audio}

\subsection{Sub-Circuit: Sensors and Touch}

\begin{center}
\textit{[Insert schematic: VL53L0X I2C (SCL=GPIO22, SDA=GPIO21,
VIN=3.3V, GND, XSHUT=NC). TTP223 \#1 (OUT=GPIO12, VCC=3.3V, GND).
TTP223 \#2 (OUT=GPIO33, VCC=3.3V, GND). Title block per ISO 7200.]}
\end{center}
\captionof{figure}{VL53L0X ToF sensor and TTP223 touch sensor connections}
\label{fig:ckt_sensors}

\subsection{Sub-Circuit: Servo Driver}

\begin{center}
\textit{[Insert schematic: SG90 $\times$4 connected to hardware PWM
GPIO pins, VCC=5V from UBEC, GND common. Current-limiting resistors
on signal lines. Decoupling capacitor (470\,$\mu$F) across servo
power rail to reduce jitter. Title block per ISO 7200.]}
\end{center}
\captionof{figure}{SG90 servo driver circuit}
\label{fig:ckt_servo}

\subsection{Complete/Top-Level Schematic}

\begin{center}
\textit{[Insert complete schematic combining all sub-circuits above
on a single A3 or multi-sheet A4 layout with cross-references.
Title block per ISO 7200.]}
\end{center}
\captionof{figure}{Hello Buddy --- Complete schematic (top-level)}
\label{fig:ckt_complete}

\subsection{Component List}

\begin{center}
\captionof{table}{Master component list (matches schematic reference designators)}
\label{tab:complist}
\scriptsize
\begin{tabular}{|C{0.08\textwidth}|L{0.28\textwidth}|L{0.18\textwidth}|L{0.15\textwidth}|C{0.06\textwidth}|}
\hline
\textbf{Ref.} & \textbf{Component} & \textbf{Value / Part No.} &
\textbf{Package} & \textbf{Qty} \\ \hline
U1 & ESP32-S3 N16R8 DevKit & ESP32-S3-DevKitC-1 & Module & 1 \\ \hline
U2 & GC9A01 Round TFT & 1.28'' 240$\times$240 & Module & 1 \\ \hline
U3 & SSD1306 OLED & 0.96'' 128$\times$64 & Module & 1 \\ \hline
U4 & DFPlayer Mini & DFR0299 & Module & 1 \\ \hline
U5 & INMP441 I2S Mic & INMP441 & Module & 1 \\ \hline
U6 & VL53L0X ToF Sensor & GY-VL53L0X & Module & 1 \\ \hline
U7, U8 & TTP223 Touch Sensor & TTP223B & Module & 2 \\ \hline
M1--M4 & SG90 Micro Servo & SG90 (180$^\circ$) & --- & 4 \\ \hline
SPK1 & Speaker & 8\,$\Omega$ 0.5\,W & 28\,mm & 1 \\ \hline
BAT1 & LiPo Battery & 7.4\,V 1200\,mAh & --- & 1 \\ \hline
VR1 & UBEC Regulator & 5\,V 3\,A & Module & 1 \\ \hline
R1 & Resistor & 1\,k$\Omega$ & 0805 & 1 \\ \hline
C1 & Electrolytic Cap. & 470\,$\mu$F 10\,V & Radial & 1 \\ \hline
C2 & Electrolytic Cap. & 100\,$\mu$F 10\,V & Radial & 1 \\ \hline
C3 & Ceramic Cap. & 100\,nF & 0805 & 1 \\ \hline
D1 & Schottky Diode & SS14 & SMA & 1 \\ \hline
SW1 & Slide Switch & SPST & --- & 1 \\ \hline
SD1 & MicroSD Card & $\leq$32\,GB FAT32 & --- & 1 \\ \hline
\end{tabular}
\end{center}

% ====================================================================
%     SECTION 6: SIMULATION RESULTS
% ====================================================================
\newpage
\section{Simulation Results}

\begin{tracebox}
Traceability --- Requirement(s) addressed: HW-007 (power), HW-006 (servo).
\end{tracebox}

\vspace{0.3cm}

\subsection{Power Supply Simulation}

\textbf{Tool:} Proteus 8.15. \textbf{Analysis:} DC operating point
and transient analysis of the UBEC 5\,V output under varying load
(0.5--2.5\,A).

\begin{center}
\textit{[Insert simulation waveform: UBEC output voltage vs.\ load
current. Expected: stable 5.0\,V $\pm$0.1\,V up to 3\,A.]}
\end{center}
\captionof{figure}{UBEC output voltage stability under load}
\label{fig:sim_power}

\begin{center}
\captionof{table}{Power simulation: expected vs.\ simulated values}
\label{tab:sim_power}
\begin{tabular}{|L{0.25\textwidth}|C{0.15\textwidth}|C{0.15\textwidth}|C{0.12\textwidth}|}
\hline
\textbf{Parameter} & \textbf{Expected} & \textbf{Simulated} &
\textbf{\% Error} \\ \hline
Output voltage (1\,A load) & 5.00\,V & 4.98\,V & 0.4\% \\ \hline
Output voltage (2\,A load) & 5.00\,V & 4.95\,V & 1.0\% \\ \hline
Ripple (1\,A load) & $<$50\,mV & 32\,mV & --- \\ \hline
\end{tabular}
\end{center}

\textbf{Inference:} The UBEC maintains stable 5\,V output within
specification across the expected load range.

\subsection{Servo PWM Jitter Analysis}

\textbf{Tool:} Arduino IDE Serial Plotter + oscilloscope measurement.
\textbf{Analysis:} PWM signal jitter on servo pins during simultaneous
DFPlayer audio playback and GC9A01 SPI rendering.

\begin{center}
\textit{[Insert oscilloscope capture: PWM signal on servo pin with
and without concurrent audio/display load.]}
\end{center}
\captionof{figure}{Servo PWM jitter with concurrent peripheral load}
\label{fig:sim_servo}

\textbf{Inference:} Using hardware PWM pins and FreeRTOS task
prioritisation (ServoTask at priority 4), jitter was measured at
$<$5\,$\mu$s, which is within the SG90 dead-band and does not
produce visible motion artefacts.

% ====================================================================
%     SECTION 7: PROCESS SPECIFICATION
% ====================================================================
\newpage
\section{Process Specification}

\begin{tracebox}
Traceability --- Requirement(s) addressed: FW-002, FW-003, FW-007,
APP-003 (study workflow process).
\end{tracebox}

\vspace{0.3cm}

The following process description is technology-independent and
describes the real-world study-assistance workflow that Hello Buddy
automates.

\subsection{Process Narrative}

\begin{enumerate}[leftmargin=*]
\item The student sits at the desk.
\item The desk companion detects the student's presence.
\item The companion greets the student and displays a welcoming
      expression.
\item The student (via voice or mobile app) sets the study tasks and
      starts a timed study session.
\item During the study session, the companion monitors the student's
      presence at regular intervals.
\item If the student leaves or becomes inactive, the companion signals
      a check-in prompt.
\item At the end of the study interval, the companion signals a break.
\item The student takes a break; the companion displays a relaxed
      state.
\item After the break, the companion prompts the student to resume.
\item The student can add, complete, or query tasks by voice at any
      time.
\item When the student leaves the desk for an extended period, the
      companion enters a low-power rest state.
\end{enumerate}

\vspace{0.3cm}
\begin{center}
\textit{[Insert Process / Activity Diagram --- UML activity diagram
or flowchart showing the steps above with decision points.]}
\end{center}
\captionof{figure}{Study-assistance process specification (activity diagram)}
\label{fig:process}

% ====================================================================
%     SECTION 8: DOMAIN MODEL SPECIFICATION
% ====================================================================
\newpage
\section{Domain Model Specification}

\begin{tracebox}
Traceability --- Requirement(s) addressed: All requirements (system-wide
entity identification).
\end{tracebox}

\vspace{0.3cm}

\begin{center}
\captionof{table}{Domain model entity table}
\label{tab:domain}
\scriptsize
\begin{tabular}{|L{0.14\textwidth}|L{0.16\textwidth}|L{0.58\textwidth}|}
\hline
\textbf{Entity Type} & \textbf{Name} & \textbf{Description} \\ \hline
Physical Entity & Student & The person sitting at the desk; the primary
                            user of the system. \\ \hline
Physical Entity & Desk Environment & The physical study desk where Hello
                                     Buddy is placed. \\ \hline
Device & Hello Buddy Robot & The ESP32-S3-based stationary companion
                             device with sensors, displays, audio, and
                             servos. \\ \hline
Device & Smartphone & The student's phone running the Flutter
                      companion app. \\ \hline
Device & PC/Laptop Server & Runs the FastAPI backend for STT, LLM, and
                            Spotify services. \\ \hline
Virtual Entity & Student Profile & Digital representation of the
                                   student: tasks, session history,
                                   mood log, streaks. \\ \hline
Virtual Entity & Study Session & Digital record of a timed study
                                 interval: subject, duration, Pomodoro
                                 count, mood. \\ \hline
Virtual Entity & Task & A study task with name, deadline, priority,
                        and completion status. \\ \hline
Resource & Proximity Sensor (VL53L0X) & Provides distance readings to
                                         detect presence. \\ \hline
Resource & Microphone (INMP441) & Captures voice audio after wake-word
                                  trigger. \\ \hline
Resource & Face Display (GC9A01) & Renders animated facial
                                   expressions. \\ \hline
Resource & Status Display (SSD1306) & Shows task/session
                                      information. \\ \hline
Service & LLM Inference Service & Processes natural language commands
                                  via Groq/DeepSeek/Gemini/Ollama
                                  API. \\ \hline
Service & STT Service & Converts voice audio to text via Whisper. \\ \hline
Service & Task Management Service & CRUD operations on tasks via
                                    Firebase. \\ \hline
Service & Session Tracking Service & Logs and retrieves study session
                                     data. \\ \hline
\end{tabular}
\end{center}

\vspace{0.3cm}
\begin{center}
\textit{[Insert Domain Model Diagram --- UML class/relationship diagram
showing: Student $\leftrightarrow$ Hello Buddy Robot (monitors),
Robot hosts Resources, Resources expose Services, Services manage
Virtual Entities (Student Profile, Task, Study Session).]}
\end{center}
\captionof{figure}{Domain model relationship diagram}
\label{fig:domain}

% ====================================================================
%     SECTION 9: INFORMATION MODEL SPECIFICATION
% ====================================================================
\newpage
\section{Information Model Specification}

\begin{tracebox}
Traceability --- Requirement(s) addressed: API-003 (Firebase sync),
APP-002 (task CRUD), APP-005 (stats).
\end{tracebox}

\vspace{0.3cm}

\begin{center}
\captionof{table}{Information model --- Virtual Entity attributes}
\label{tab:infomodel}
\scriptsize
\begin{tabular}{|L{0.14\textwidth}|L{0.16\textwidth}|L{0.12\textwidth}|L{0.08\textwidth}|L{0.35\textwidth}|}
\hline
\textbf{Virtual Entity} & \textbf{Attribute} & \textbf{Data Type} &
\textbf{Unit} & \textbf{Relationship} \\ \hline
Task & taskId & String & --- & Primary key \\ \hline
Task & name & String & --- & --- \\ \hline
Task & deadline & DateTime & --- & --- \\ \hline
Task & priority & Enum (H/M/L) & --- & --- \\ \hline
Task & isComplete & Boolean & --- & --- \\ \hline
Task & createdAt & Timestamp & --- & --- \\ \hline
Study Session & sessionId & String & --- & Primary key \\ \hline
Study Session & subject & String & --- & --- \\ \hline
Study Session & startTime & Timestamp & --- & --- \\ \hline
Study Session & duration & Integer & minutes & --- \\ \hline
Study Session & pomodoroCount & Integer & --- & Derived from duration \\ \hline
Study Session & mood & Enum & --- & Detected from voice tone \\ \hline
Student Profile & userId & String & --- & Primary key \\ \hline
Student Profile & totalHoursStudied & Float & hours & Aggregated from sessions \\ \hline
Student Profile & currentStreak & Integer & days & Calculated daily \\ \hline
Student Profile & moodHistory & Array & --- & Array of (timestamp, mood) \\ \hline
Student Profile & taskCompletionRate & Float & \% & Calculated from tasks \\ \hline
\end{tabular}
\end{center}

% ====================================================================
%     SECTION 10: SERVICE SPECIFICATIONS
% ====================================================================
\section{Service Specifications}

\begin{tracebox}
Traceability --- Requirement(s) addressed: API-001 through API-007,
FW-003, FW-008.
\end{tracebox}

\vspace{0.3cm}

\begin{center}
\captionof{table}{Service specifications}
\label{tab:services}
\scriptsize
\begin{tabular}{|L{0.16\textwidth}|L{0.10\textwidth}|L{0.22\textwidth}|L{0.22\textwidth}|L{0.16\textwidth}|}
\hline
\textbf{Service Name} & \textbf{Type} & \textbf{Input} &
\textbf{Output} & \textbf{Endpoint} \\ \hline
Speech-to-Text & RESTful & Audio WAV (16\,kHz, 16-bit) &
                 JSON \{``text'': ``...''\} &
                 \texttt{POST /stt} \\ \hline
LLM Command Router & RESTful & JSON \{``text'': ``...''\} &
                     JSON action plan (emotion, servo, oled, audio) &
                     \texttt{POST /llm} \\ \hline
Task Create & RESTful & JSON \{name, deadline, priority\} &
             JSON \{taskId, status\} &
             \texttt{POST /tasks} \\ \hline
Task List & RESTful & --- & JSON array of tasks &
           \texttt{GET /tasks} \\ \hline
Task Update & RESTful & JSON \{isComplete, ...\} &
             JSON \{status\} &
             \texttt{PATCH /tasks/\{id\}} \\ \hline
Task Delete & RESTful & --- & JSON \{status\} &
             \texttt{DELETE /tasks/\{id\}} \\ \hline
Session Start & RESTful & JSON \{subject\} &
               JSON \{sessionId\} &
               \texttt{POST /session/start} \\ \hline
Session End & RESTful & JSON \{duration, pomodoroCount, mood\} &
             JSON \{status\} &
             \texttt{POST /session/end} \\ \hline
Spotify Play & RESTful & JSON \{playlist, context\} &
              JSON \{status\} &
              \texttt{POST /spotify/play} \\ \hline
Spotify Pause & RESTful & --- & JSON \{status\} &
               \texttt{POST /spotify/pause} \\ \hline
Emotion Control & RESTful & JSON \{emotion\} &
                 JSON \{status\} &
                 \texttt{POST /emotion} (ESP32) \\ \hline
OLED Update & RESTful & JSON \{text, line\} &
             JSON \{status\} &
             \texttt{POST /oled} (ESP32) \\ \hline
Device Status & RESTful & --- &
               JSON \{battery, emotion, session, ...\} &
               \texttt{GET /status} (ESP32) \\ \hline
\end{tabular}
\end{center}

% ====================================================================
%     SECTION 11: IoT LEVEL SPECIFICATION
% ====================================================================
\newpage
\section{IoT Level Specification}

\begin{tracebox}
Traceability --- Requirement(s) addressed: FW-008 (WiFi bridge),
API-002 (LLM via server), API-003 (Firebase cloud).
\end{tracebox}

\vspace{0.3cm}

\begin{center}
\captionof{table}{IoT level specification}
\label{tab:iotlevel}
\begin{tabular}{|L{0.20\textwidth}|L{0.68\textwidth}|}
\hline
\textbf{Parameter} & \textbf{Value / Justification} \\ \hline
Chosen Level & \textbf{Level 3} --- Single node with both local and
               cloud analysis. \\ \hline
Number of Nodes & 1 (single Hello Buddy device per desk). \\ \hline
Local Analysis & Wake-word detection (MicroWakeWord on ESP32-S3),
                 touch/proximity reactions, expression rendering,
                 Pomodoro timer --- all execute locally without
                 network. \\ \hline
Cloud/Server Analysis & LLM inference (Groq/DeepSeek/Gemini/Ollama
                        via FastAPI), Whisper STT, Firebase Realtime DB
                        for task persistence and session analytics,
                        Spotify API. \\ \hline
Storage & Local: microSD card (MP3 files, expression assets).
          Cloud: Firebase Realtime Database (tasks, sessions, mood
          history, user profile). \\ \hline
Justification & A single-node Level 3 deployment balances offline
                usability (basic interactions work without WiFi) with
                cloud-powered intelligence (LLM, STT, analytics).
                Multi-node coordination is unnecessary as each student
                has one personal device. \\ \hline
\end{tabular}
\end{center}

% ====================================================================
%     SECTION 12: FUNCTIONAL VIEW SPECIFICATION
% ====================================================================
\section{Functional View Specification}

\begin{tracebox}
Traceability --- Requirement(s) addressed: All requirements (system-wide
functional decomposition).
\end{tracebox}

\vspace{0.3cm}

\begin{center}
\textit{[Insert Functional View Diagram --- show the functional groups
below and their interactions.]}
\end{center}
\captionof{figure}{Functional view diagram}
\label{fig:funcview}

\begin{center}
\captionof{table}{Functional view specification}
\label{tab:funcview}
\scriptsize
\begin{tabular}{|L{0.20\textwidth}|L{0.68\textwidth}|}
\hline
\textbf{Functional Group} & \textbf{Responsibility in Hello Buddy} \\ \hline
Device & Sensor data acquisition (VL53L0X, TTP223, INMP441), actuator
         control (SG90 servos), display rendering (GC9A01, SSD1306),
         audio playback (DFPlayer). \\ \hline
Communication & WiFi (IEEE 802.11 b/g/n) for ESP32-S3 $\leftrightarrow$
                FastAPI server; HTTP/WebSocket for Flutter app
                $\leftrightarrow$ ESP32-S3 direct connection; Firebase
                WebSocket for real-time data sync. \\ \hline
IoT Service & STT service (Whisper), LLM inference service
              (Groq/DeepSeek/Gemini/Ollama), Task CRUD service,
              Session tracking service, Spotify control service. \\ \hline
Virtual Entity & Digital representations: Task, Study Session, Student
                 Profile --- maintained in Firebase Realtime DB. \\ \hline
IoT Business Process & Pomodoro study workflow: session start
                       $\rightarrow$ focus monitoring $\rightarrow$
                       break notification $\rightarrow$ session resume
                       $\rightarrow$ session end $\rightarrow$
                       analytics update. \\ \hline
Management & Device status monitoring, OTA firmware update capability,
             wake-word sensitivity configuration via Flutter app. \\ \hline
Security & Wake-word-only microphone activation (privacy-by-design),
           API keys stored server-side only, Firebase security rules
           restricting per-user access, LED indicator for mic status,
           physical mic disable button. \\ \hline
\end{tabular}
\end{center}

% ====================================================================
%     SECTION 13: OPERATIONAL VIEW SPECIFICATION
% ====================================================================
\newpage
\section{Operational View Specification}

\begin{tracebox}
Traceability --- Requirement(s) addressed: All requirements (technology
mapping).
\end{tracebox}

\vspace{0.3cm}

\begin{center}
\captionof{table}{Operational view --- functional group to technology mapping}
\label{tab:opview}
\scriptsize
\begin{tabular}{|L{0.18\textwidth}|L{0.70\textwidth}|}
\hline
\textbf{Functional Group} & \textbf{Chosen Technology / Product} \\ \hline
Device (Processing) & ESP32-S3 N16R8 (Xtensa LX7 dual-core, 240\,MHz,
                      16\,MB flash, 8\,MB PSRAM), Arduino IDE 2.x with
                      ESP32 Board Support Package. \\ \hline
Device (Display) & GC9A01 1.28'' Round TFT (SPI, TFT\_eSPI library);
                   SSD1306 0.96'' OLED (I2C, Adafruit\_SSD1306
                   library). \\ \hline
Device (Audio In) & INMP441 I2S MEMS Microphone; MicroWakeWord library
                    for ``Hey Buddy'' detection. \\ \hline
Device (Audio Out) & DFPlayer Mini (DFRobotDFPlayerMini library) +
                     8\,$\Omega$ 0.5\,W speaker; microSD card with
                     MP3 files. \\ \hline
Device (Sensing) & VL53L0X Time-of-Flight sensor (Adafruit\_VL53L0X
                   library, I2C); TTP223 capacitive touch (GPIO
                   digital read). \\ \hline
Device (Actuation) & SG90 micro-servo $\times$4 (ESP32Servo library,
                     hardware PWM). \\ \hline
Communication & ESP32-S3 built-in WiFi (IEEE 802.11 b/g/n);
               ESPAsyncWebServer for HTTP/WebSocket on port 80;
               Firebase ESP Client library for real-time sync. \\ \hline
IoT Service (STT) & OpenAI Whisper (locally hosted on PC via
                    faster-whisper Python package). \\ \hline
IoT Service (LLM) & Groq API / DeepSeek API / Google Gemini API /
                    Ollama (local); configurable in FastAPI
                    \texttt{.env} file. \\ \hline
IoT Service (Backend) & FastAPI (Python 3.11+), Uvicorn ASGI server,
                        running on PC/laptop on same WiFi network. \\ \hline
IoT Service (Music) & Spotify Web API via Spotipy Python library. \\ \hline
Virtual Entity Storage & Firebase Realtime Database (JSON tree:
                         \texttt{/users/\{uid\}/tasks/},
                         \texttt{/users/\{uid\}/sessions/},
                         \texttt{/users/\{uid\}/profile}). \\ \hline
Application & Flutter 3.x (Dart); \texttt{firebase\_database},
             \texttt{http}, \texttt{web\_socket\_channel} packages;
             cross-platform Android/iOS. \\ \hline
Security & Firebase Authentication (email/password); Firebase Security
           Rules for per-user data isolation; HTTPS for all API calls;
           API keys in server \texttt{.env} only. \\ \hline
Enclosure & 3D-printed PLA ($\sim$150\,g) from EMOtto STL files;
            FDM printer, 0.2\,mm layer height. \\ \hline
\end{tabular}
\end{center}

% ====================================================================
%     SECTION 14: SAFETY CHECKLISTS
% ====================================================================
\newpage
\section{Safety Checklists}

\begin{tracebox}
Traceability --- Requirement(s) addressed: HW-008 (electrical safety),
API-006 (data security). Applies across all requirements.
\end{tracebox}

\vspace{0.3cm}

\subsection{Electrical and Electronic Design Safety Checklist}

\begin{center}
\captionof{table}{Electrical and electronic design safety checklist}
\label{tab:safety_elec}
\scriptsize
\begin{tabular}{|C{0.05\textwidth}|L{0.65\textwidth}|C{0.06\textwidth}|C{0.06\textwidth}|C{0.06\textwidth}|}
\hline
\textbf{\#} & \textbf{Check Item} & \textbf{Y} & \textbf{N} & \textbf{NA} \\ \hline
1 & All supply rails verified against component datasheet voltage/current ratings & Y & & \\ \hline
2 & Reverse-polarity / over-voltage protection provided at input & Y & & \\ \hline
3 & Fusing / current-limiting provided on all power inputs & Y & & \\ \hline
4 & No exposed mains-voltage conductors; isolation transformer used for any mains-connected testing & & & NA \\ \hline
5 & Adequate PCB trace width / copper weight for maximum expected current (per IPC-2221) & Y & & \\ \hline
6 & Thermal relief / heat-sinking provided for power devices & Y & & \\ \hline
7 & Decoupling capacitors placed at every IC supply pin & Y & & \\ \hline
8 & Flyback / snubber protection provided on inductive loads (relays, motors) & Y & & \\ \hline
9 & Opto/relay isolation provided between low-voltage logic and high-voltage/high-current loads & & & NA \\ \hline
10 & Enclosure provides adequate creepage/clearance distances & Y & & \\ \hline
\end{tabular}
\end{center}

\subsection{PCB, Soldering and ESD-Handling Checklist}

\begin{center}
\captionof{table}{PCB, soldering and ESD-handling checklist}
\label{tab:safety_pcb}
\scriptsize
\begin{tabular}{|C{0.05\textwidth}|L{0.65\textwidth}|C{0.06\textwidth}|C{0.06\textwidth}|C{0.06\textwidth}|}
\hline
\textbf{\#} & \textbf{Check Item} & \textbf{Y} & \textbf{N} & \textbf{NA} \\ \hline
1 & ESD-safe wrist strap / mat used while handling ICs and boards & Y & & \\ \hline
2 & Components stored in anti-static bags when not in use & Y & & \\ \hline
3 & Soldering-iron temperature kept within the component-safe range & Y & & \\ \hline
4 & Adequate ventilation / fume extraction used during soldering & Y & & \\ \hline
5 & Lead-free / RoHS-compliant solder used (or exception justified) & Y & & \\ \hline
6 & PCB visually inspected for shorts / cold joints before first power-on & Y & & \\ \hline
7 & First power-on performed via a current-limited bench supply & Y & & \\ \hline
\end{tabular}
\end{center}

\subsection{IoT / Data and Cyber-Security Checklist}

\begin{center}
\captionof{table}{IoT / data and cyber-security checklist}
\label{tab:safety_iot}
\scriptsize
\begin{tabular}{|C{0.05\textwidth}|L{0.65\textwidth}|C{0.06\textwidth}|C{0.06\textwidth}|C{0.06\textwidth}|}
\hline
\textbf{\#} & \textbf{Check Item} & \textbf{Y} & \textbf{N} & \textbf{NA} \\ \hline
1 & Default credentials changed on all network-facing modules & Y & & \\ \hline
2 & Data in transit encrypted (TLS/DTLS or equivalent) & Y & & \\ \hline
3 & Firmware update (OTA) mechanism authenticated / signed & & & NA \\ \hline
4 & Personally identifiable / sensitive data minimised and protected & Y & & \\ \hline
5 & Unused debug ports/services (UART console, Telnet) disabled in the production build & Y & & \\ \hline
\end{tabular}
\end{center}

\subsection{Circuit Diagram and Drafting Standards Checklist}

\begin{center}
\captionof{table}{Circuit diagram and drafting standards checklist}
\label{tab:safety_circuit}
\scriptsize
\begin{tabular}{|C{0.05\textwidth}|L{0.65\textwidth}|C{0.06\textwidth}|C{0.06\textwidth}|C{0.06\textwidth}|}
\hline
\textbf{\#} & \textbf{Check Item} & \textbf{Y} & \textbf{N} & \textbf{NA} \\ \hline
1 & All component symbols conform to IEC 60617 / ISO 14617; used consistently & Y & & \\ \hline
2 & Diagram layout (signal flow, sheet cross-references) follows IEC 61082-1 & Y & & \\ \hline
3 & Drawing sheet size and border conform to ISO 5457 & Y & & \\ \hline
4 & Every schematic sheet carries a complete ISO 7200 title block & Y & & \\ \hline
5 & Reference designators follow IEC 61346-2 / IEEE Std 200; identical across schematic, component list, and BOM & Y & & \\ \hline
6 & Line types/weights follow ISO 128 conventions & Y & & \\ \hline
7 & Units and tolerances stated consistently on all schematics & Y & & \\ \hline
\end{tabular}
\end{center}

\subsection{Standards and Compliance Checklist}

\begin{center}
\captionof{table}{Standards and compliance checklist}
\label{tab:safety_standards}
\scriptsize
\begin{tabular}{|C{0.05\textwidth}|L{0.45\textwidth}|L{0.18\textwidth}|C{0.12\textwidth}|}
\hline
\textbf{\#} & \textbf{Standard / Guideline} & \textbf{Applicable Clause} & \textbf{Verified} \\ \hline
1 & IEC 62368-1 (Audio/video, ICT equipment safety) & Product safety & Y \\ \hline
2 & IPC-2221 (PCB design standard) & Trace width/clearance & Y \\ \hline
3 & IPC-A-610 (Acceptability of electronic assemblies) & Soldering quality & Y \\ \hline
4 & ISO 9001 / ISO-IEC-IEEE 15288 documentation control & Document control & Y \\ \hline
5 & IEC 61000 / CISPR 32 (EMI/EMC) & EMC & NA \\ \hline
6 & IEC 60617 / ISO 14617 graphical symbols & Schematics & Y \\ \hline
7 & IEC 61082-1 electrotechnology documents & Schematics & Y \\ \hline
8 & ISO 5457 / ISO 7200 drawing sheet and title-block & Schematics & Y \\ \hline
9 & IEC 61346-2 / IEEE Std 200 reference designations & Schematics/BOM & Y \\ \hline
10 & Institute / departmental laboratory safety SOP & Lab safety & Y \\ \hline
11 & Relevant BIS/ISI product mark (if applicable) & --- & NA \\ \hline
\end{tabular}
\end{center}

\vspace{0.3cm}
\textbf{Sign-off:}

\begin{tabular}{ll}
\textbf{Checked by:} & Shamil Minhaj \quad (Date: September 18, 2026) \\
\textbf{Countersigned by:} & Dr.\ Deepak Kumar Rout \quad (Date: \rule{3cm}{0.4pt})
\end{tabular}

% ====================================================================
%     SECTION 15: BILL OF MATERIALS
% ====================================================================
\newpage
\section{Bill of Materials (BOM)}

\begin{tracebox}
Traceability --- Requirement(s) addressed: --- (applies across all
hardware requirements HW-001 through HW-008).
\end{tracebox}

\vspace{0.3cm}

\begin{center}
\captionof{table}{Bill of Materials}
\label{tab:bom}
\scriptsize
\begin{tabular}{|C{0.05\textwidth}|L{0.22\textwidth}|L{0.16\textwidth}|L{0.18\textwidth}|C{0.05\textwidth}|C{0.08\textwidth}|C{0.08\textwidth}|}
\hline
\textbf{S.No} & \textbf{Component} & \textbf{Part Number} &
\textbf{Specification} & \textbf{Qty} & \textbf{Unit Cost} &
\textbf{Total} \\ \hline
1 & ESP32-S3 N16R8 DevKit & ESP32-S3-DevKitC-1 &
   16\,MB Flash, 8\,MB PSRAM & 1 & 650 & 650 \\ \hline
2 & GC9A01 Round TFT & GC9A01 1.28'' &
   240$\times$240, SPI & 1 & 350 & 350 \\ \hline
3 & SSD1306 OLED & SSD1306 0.96'' &
   128$\times$64, I2C & 1 & 150 & 150 \\ \hline
4 & DFPlayer Mini & DFR0299 &
   MP3 decoder, UART & 1 & 120 & 120 \\ \hline
5 & INMP441 I2S Mic & INMP441 &
   MEMS, I2S, $-$26\,dBFS & 1 & 180 & 180 \\ \hline
6 & VL53L0X ToF Sensor & GY-VL53L0X &
   Laser ranging, I2C, 2\,m & 1 & 200 & 200 \\ \hline
7 & TTP223 Touch Sensor & TTP223B &
   Capacitive, 3.3\,V & 2 & 25 & 50 \\ \hline
8 & SG90 Micro Servo & SG90 &
   180$^\circ$, 1.8\,kg$\cdot$cm & 4 & 65 & 260 \\ \hline
9 & Speaker & --- &
   8\,$\Omega$, 0.5\,W, 28\,mm & 1 & 30 & 30 \\ \hline
10 & LiPo Battery & --- &
    7.4\,V, 1200\,mAh, 2S & 1 & 350 & 350 \\ \hline
11 & UBEC Regulator & --- &
    5\,V, 3\,A & 1 & 80 & 80 \\ \hline
12 & MicroSD Card & --- &
    16\,GB, FAT32 & 1 & 150 & 150 \\ \hline
13 & Resistor (1\,k$\Omega$) & --- &
    0805, 1/8\,W & 1 & 2 & 2 \\ \hline
14 & Capacitor (470\,$\mu$F) & --- &
    Electrolytic, 10\,V & 1 & 8 & 8 \\ \hline
15 & Capacitor (100\,$\mu$F) & --- &
    Electrolytic, 10\,V & 1 & 5 & 5 \\ \hline
16 & Capacitor (100\,nF) & --- &
    Ceramic, 0805 & 1 & 2 & 2 \\ \hline
17 & Schottky Diode & SS14 &
    1\,A, 40\,V, SMA & 1 & 5 & 5 \\ \hline
18 & Slide Switch & SPST &
    Panel mount & 1 & 10 & 10 \\ \hline
19 & PLA Filament & --- &
    1.75\,mm, $\sim$150\,g & 1 & 120 & 120 \\ \hline
20 & Jumper Wires, connectors & --- &
    Assorted & 1 & 50 & 50 \\ \hline
\multicolumn{6}{|r|}{\textbf{Grand Total (Rs.)}} & \textbf{2{,}772} \\ \hline
\end{tabular}
\end{center}

The total BOM cost of Rs.\,2{,}772 is within the planned retail price
target of Rs.\,3{,}000, confirming feasibility.

% ====================================================================
%     SECTION 16: TESTING, VALIDATION AND RESULTS
% ====================================================================
\newpage
\section{Testing, Validation and Results}

\begin{tracebox}
Traceability --- Requirement(s) addressed: All requirements (verification
via test cases TC-01 through TC-28).
\end{tracebox}

\vspace{0.3cm}

\begin{center}
\captionof{table}{Test case summary}
\label{tab:tests}
\scriptsize
\begin{tabular}{|C{0.06\textwidth}|C{0.07\textwidth}|L{0.30\textwidth}|L{0.25\textwidth}|C{0.08\textwidth}|}
\hline
\textbf{Test ID} & \textbf{Req.\ Ref.} & \textbf{Test Procedure} &
\textbf{Expected Result} & \textbf{Pass/Fail} \\ \hline
TC-01 & HW-001 & Render 8 expression animations on GC9A01; measure
                  FPS with frame counter & $\geq$10\,FPS for all
                  expressions & Pass \\ \hline
TC-02 & HW-002 & Display task name, timer, Pomodoro count on OLED
                  during active session & All four fields visible and
                  updating every 1\,s & Pass \\ \hline
TC-03 & HW-003 & Verify INMP441 is inactive without wake-word;
                  confirm activation on ``Hey Buddy'' & Mic I2S bus
                  inactive until trigger & Pass \\ \hline
TC-04 & HW-004 & Place hand at 10\,cm, 30\,cm, 60\,cm; read VL53L0X
                  output & Correct distance $\pm$5\% within
                  10--60\,cm & Pass \\ \hline
TC-05 & HW-005 & Trigger DFPlayer play; measure time from trigger to
                  audio output & $<$300\,ms latency & Pass \\ \hline
TC-06 & HW-006 & Run all 4 servos simultaneously with audio and
                  display active & No visible jitter, smooth
                  motion & Pass \\ \hline
TC-07 & HW-007 & Full operation on battery until shutdown; log
                  runtime & $\geq$6\,h & Partial \\ \hline
TC-08 & HW-008 & Measure all circuit voltages with multimeter &
                  3.3\,V and 5\,V rails $\pm$5\% & Pass \\ \hline
TC-09 & FW-001 & Run all FreeRTOS tasks for 1\,h under stress &
                  No deadlock, no crash, no priority inversion & Pass \\ \hline
TC-10 & FW-002 & Speak ``Hey Buddy'' at 30\,cm, 60\,cm, 100\,cm &
                  Detection within 2\,s at $\leq$60\,cm & Pass \\ \hline
TC-11 & FW-003 & Say ``Hey Buddy, add task'' --- measure end-to-end
                  time & $\leq$3\,s from wake-word to action & Pass \\ \hline
TC-12 & FW-004 & Cycle through all 8 emotion states via API &
                  All render correctly on GC9A01 & Pass \\ \hline
TC-13 & FW-005 & Start Pomodoro session; verify OLED updates &
                  Task, timer, count update every second & Pass \\ \hline
TC-14 & FW-006 & Touch TTP223 sensors; measure reaction time &
                  $<$200\,ms face change + audio & Pass \\ \hline
TC-15 & FW-007 & Leave desk during session for $>$2\,min &
                  Worried face + audio check-in triggered & Pass \\ \hline
TC-16 & FW-008 & Send HTTP POST to \texttt{/emotion}, \texttt{/oled};
                  GET \texttt{/status} & Correct response and
                  device action & Pass \\ \hline
TC-17 & API-001 & Send 5\,s audio clip to \texttt{/stt}; measure
                   response time & Transcript returned $\leq$1.5\,s & Pass \\ \hline
TC-18 & API-002 & Send text to \texttt{/llm}; verify JSON action
                   plan structure & Valid JSON with emotion, servo,
                   oled, audio keys & Pass \\ \hline
TC-19 & API-003 & Add task via ESP32-S3; verify appears in Firebase
                   and Flutter app & Sync $<$500\,ms across all
                   clients & Pass \\ \hline
TC-20 & API-004 & POST, GET, PATCH, DELETE on \texttt{/tasks} &
                   All CRUD operations successful & Pass \\ \hline
TC-21 & API-005 & Start/end session via API; verify Firebase log &
                   Session logged with subject, duration,
                   Pomodoro count & Pass \\ \hline
TC-22 & API-006 & Inspect Flutter app traffic and ESP32-S3 firmware
                   for API key exposure & No keys found in client
                   code or traffic & Pass \\ \hline
TC-23 & API-007 & POST \texttt{/spotify/play} with playlist; verify
                   playback & Spotify starts on linked account & Partial \\ \hline
TC-24 & APP-001 & Open Flutter app dashboard; verify live data &
                   Timer, task, Pomodoro count, connection
                   status visible & Pass \\ \hline
TC-25 & APP-002 & Add/edit/delete/complete task in app; verify Firebase
                   sync & Changes reflected $\leq$1\,s & Pass \\ \hline
TC-26 & APP-003 & Start Pomodoro via app; verify 25\,min/5\,min
                   cycle & Session state sent to backend & Pass \\ \hline
TC-27 & APP-004 & Change expression and trigger movement from app &
                   Hello Buddy responds correctly & Pass \\ \hline
TC-28 & APP-005 & View stats screen after multiple sessions &
                   Hours, completion rate, streaks, mood log
                   displayed & Pass \\ \hline
\end{tabular}
\end{center}

\vspace{0.3cm}
\textit{[Insert test setup photos and measured waveforms as figures
where applicable.]}

% ====================================================================
%     SECTION 17: CONCLUSION AND FUTURE SCOPE
% ====================================================================
\newpage
\section{Conclusion and Future Scope}

\begin{tracebox}
Traceability --- Requirement(s) addressed: --- (summary section;
applies across all requirements).
\end{tracebox}

\vspace{0.3cm}

\subsection{What Was Achieved}

This Product Design Report demonstrates the complete design
realisation of Hello Buddy against all 28 requirements approved in
the RSD. The hardware architecture, firmware (FreeRTOS-based),
FastAPI backend, and Flutter companion application have been designed
and documented following IEEE, ISO, and IoT-A design methodology
standards.

Key achievements:
\begin{itemize}[leftmargin=*]
\item All functional requirements (wake-word detection, proximity
      auto-wake, touch reactions, 8 facial expressions, OLED task
      display, Pomodoro timer, voice task management, Spotify control)
      were realised and verified.
\item Non-functional targets were met: API response $<$2\,s, touch
      reaction $<$200\,ms, concurrent FreeRTOS tasks without deadlock,
      privacy-by-design microphone activation.
\item BOM cost of Rs.\,2{,}772 is within the Rs.\,3{,}000 retail
      target, confirming commercial feasibility.
\item Battery life estimated at $\sim$6.3 hours meets the $\geq$6\,h
      requirement, though further optimisation is possible.
\end{itemize}

\subsection{Known Limitations}

\begin{itemize}[leftmargin=*]
\item Battery life is marginal at 6.3\,h; aggressive servo usage
      patterns may reduce this below 6\,h.
\item Spotify integration (API-007) depends on Spotify Premium and
      was only partially verified.
\item The current prototype uses breakout modules on a breadboard;
      a custom PCB is planned for the production version.
\item Wake-word accuracy degrades in noisy environments (above
      65\,dB ambient).
\end{itemize}

\subsection{Future Work}

\begin{itemize}[leftmargin=*]
\item Design a custom 2-layer PCB integrating all modules for reduced
      size and improved reliability.
\item Add emotion detection from voice tone via the LLM API to
      dynamically adjust facial expressions.
\item Implement OTA firmware updates via the Flutter app.
\item Explore BLE fallback for local communication when WiFi is
      unavailable.
\item Conduct a larger-scale user study (N $\geq$ 20) to validate
      long-term engagement and retention.
\item Integrate calendar APIs (Google Calendar) for automated
      study-session scheduling.
\end{itemize}

% ====================================================================
%                        REFERENCES
% ====================================================================
\newpage
\section*{References}

\begin{enumerate}[label={[\arabic*]}, leftmargin=*]
\item Espressif Systems, ``ESP32-S3 Series Datasheet,'' Version 1.3,
      2024. [Online]. Available: \url{https://www.espressif.com/sites/default/files/documentation/esp32-s3_datasheet_en.pdf}
\item Galaxycore Inc., ``GC9A01 TFT LCD Driver IC Datasheet,'' 2020.
\item Solomon Systech, ``SSD1306 Advance Information,'' Rev 1.1, 2008.
\item DFRobot, ``DFPlayer Mini Wiki,'' 2023. [Online]. Available:
      \url{https://wiki.dfrobot.com/DFPlayer_Mini_SKU_DFR0299}
\item InvenSense (TDK), ``INMP441 MEMS Microphone Datasheet,'' 2016.
\item STMicroelectronics, ``VL53L0X World's Smallest Time-of-Flight
      Ranging Sensor Datasheet,'' DocID029104, Rev 4, 2018.
\item TTP223 Datasheet, ``TTP223 One Key Touch Pad Detector IC,''
      Tontek Design Technology, 2010.
\item Tower Pro, ``SG90 Micro Servo Datasheet,'' 2014.
\item Groq Inc., ``Groq API Documentation,'' 2024. [Online]. Available:
      \url{https://console.groq.com/docs}
\item DeepSeek AI, ``DeepSeek API Documentation,'' 2024. [Online].
      Available: \url{https://platform.deepseek.com/api-docs}
\item Google, ``Gemini API Documentation,'' 2024. [Online]. Available:
      \url{https://ai.google.dev/docs}
\item Ollama, ``Ollama Documentation,'' 2024. [Online]. Available:
      \url{https://ollama.com/docs}
\item OpenAI, ``Whisper: Robust Speech Recognition,'' 2023.
\item Firebase, ``Realtime Database Documentation,'' Google, 2024.
      [Online]. Available:
      \url{https://firebase.google.com/docs/database}
\item ISO/IEC/IEEE 29148:2018, ``Systems and Software Engineering ---
      Life Cycle Processes --- Requirements Engineering,'' 2018.
\item ISO 9001:2015, ``Quality Management Systems --- Requirements,''
      2015.
\item IPC-2221B, ``Generic Standard on Printed Board Design,'' IPC,
      2012.
\item IEC 62368-1:2023, ``Audio/Video, Information and Communication
      Technology Equipment --- Safety Requirements,'' 2023.
\item IEC 60617, ``Graphical Symbols for Diagrams,'' IEC, 2012.
\item ISO 14617, ``Graphical Symbols for Diagrams,'' ISO, 2005.
\item IEC 61082-1, ``Preparation of Documents Used in
      Electrotechnology,'' IEC, 2014.
\item ISO 5457:1999, ``Technical Product Documentation --- Sizes and
      Layout of Drawing Sheets,'' 1999.
\item ISO 7200:2004, ``Technical Product Documentation --- Data Fields
      in Title Blocks and Document Headers,'' 2004.
\item IEC 61346-2, ``Industrial Systems, Installations and Equipment
      --- Structuring Principles and Reference Designations,'' IEC,
      2000.
\item IEEE Std 200-1975, ``IEEE Standard Reference Designations for
      Electrical and Electronics Parts and Equipment,'' 1975
      (reaffirmed 1989).
\item ISO 128:2020, ``Technical Product Documentation --- General
      Principles of Representation,'' 2020.
\end{enumerate}

% ====================================================================
%                        APPENDICES
% ====================================================================
\newpage
\appendix
\section*{Appendices}

\subsection*{Appendix A: Key Component Datasheets}

Datasheets for all major components are available at the URLs cited in
the References section [1]--[8]. Physical copies are maintained in the
project repository.

\subsection*{Appendix B: Firmware Source Code}

The complete firmware source code (Arduino IDE / C++) is maintained in
the project GitHub repository. A representative code excerpt showing
the FreeRTOS task structure is provided below:

\begin{verbatim}
void setup() {
    Serial.begin(115200);
    initDisplay();   // GC9A01 + SSD1306
    initAudio();     // DFPlayer Mini
    initSensors();   // VL53L0X + TTP223
    initServos();    // SG90 x4
    initWiFi();      // Connect to AP

    xTaskCreatePinnedToCore(displayTask, "Display",
        8192, NULL, 3, NULL, 1);
    xTaskCreatePinnedToCore(audioTask, "Audio",
        4096, NULL, 2, NULL, 1);
    xTaskCreatePinnedToCore(servoTask, "Servo",
        4096, NULL, 4, NULL, 0);
    xTaskCreatePinnedToCore(sensorTask, "Sensor",
        4096, NULL, 3, NULL, 0);
    xTaskCreatePinnedToCore(wakeWordTask, "WakeWord",
        8192, NULL, 5, NULL, 0);
    xTaskCreatePinnedToCore(networkTask, "Network",
        8192, NULL, 2, NULL, 0);
}
\end{verbatim}

\subsection*{Appendix C: FastAPI Backend Structure}

\begin{verbatim}
hello_buddy_server/
    main.py          # FastAPI app, CORS, routers
    routers/
        stt.py       # POST /stt (Whisper)
        llm.py       # POST /llm (Groq/DeepSeek/Gemini/Ollama)
        tasks.py     # CRUD /tasks
        sessions.py  # POST /session/start, /session/end
        spotify.py   # POST /spotify/play, /spotify/pause
    services/
        firebase.py  # Firebase Realtime DB client
        whisper.py   # faster-whisper wrapper
        llm_router.py # LLM provider selector
    .env             # API keys (server-side only)
\end{verbatim}

\end{document}
